\section{Implementation Details}
\label{app:details}

\subsection{State schema and controller invariants}

Every graph object has a stable ID, type, creation step, provenance pointer, and status. Evidence nodes additionally record corpus document and span IDs; claims record premise IDs and the operation that produced them; answer candidates retain their own support channels. Hyperedges are directed and typed. A branch is a view over active nodes and edges, not an independent mutable graph.

The controller validates the following invariants before and after every committed transition:
\begin{enumerate}[leftmargin=*]
  \item only the controller mutates graph, budget, branch, or terminal state;
  \item every derived claim has valid and earlier provenance;
  \item active hyperedges reference extant, type-compatible premises and conclusions;
  \item the budget ledger is monotone and reconciles predicted with actual resource use;
  \item terminal states are mutually exclusive and immutable;
  \item evaluation-only annotations and gold answers are unreachable from runtime prompts, retrieval queries, scoring, and control flow.
\end{enumerate}
A failed validation rejects the proposal and emits an auditable error event rather than partially mutating the state.

\subsection{Belief channels}

Absolute support $s_v$ estimates whether node $v$ is independently supported. Relative weight $w_v$ compares mutually exclusive candidates but is never used as a substitute for $s_v$. Entropy $h_v$ summarizes disagreement among eligible alternatives. Evidence gap $g_v$ measures missing support required by the dependency schema. Contradiction mass $k_v$ aggregates active counterevidence. Downstream impact $d_v$ is computed from reachable answer dependencies, while dependency unlock $u_v$ estimates how many blocked joins or terminal gates would become feasible if $v$ were resolved. Heat $z_v$ is the typed propagated demand in Equation~\ref{eq:diffusion}.

All channels are stored and logged separately. Scalar packet scores are ephemeral controller decisions; they do not overwrite the underlying belief vector.

\subsection{Operation--fidelity packets}

\begin{table}[h]
\caption{Packet families and illustrative fidelity controls. Exact production settings remain frozen in configuration files.}
\label{tab:packets}
\centering
\begin{tabular}{lll}
\toprule
Operation & Target & Fidelity controls \\
\midrule
\textsc{Retrieve} & subgoal/claim & query count, top-$k$, reranking \\
\textsc{Verify} & claim/candidate & context size, independent samples \\
\textsc{Join} & premise set & schema breadth, verifier depth \\
\textsc{Revise} & conflict region & local/global edit scope \\
\textsc{Consolidate} & duplicates & merge breadth, provenance checks \\
\textsc{AnswerCheck} & candidate & support and contradiction checks \\
\bottomrule
\end{tabular}
\end{table}

Packet feasibility is determined before scoring. Infeasible packets include operations with missing typed inputs, repeated no-op signatures, exceeded per-family limits, or predicted costs outside the residual budget. Stable tie-breaking uses packet family, target ID, and fidelity ID, making controller behavior deterministic conditional on model and retriever outputs.

\subsection{Actual-utility feedback}

For an executed packet $p_t$, realized utility is a bounded, normalized change in graph readiness:
\begin{align}
r_t={}&\beta_1\Delta \overline{s}
-\beta_2\Delta \overline{g}
-\beta_3\Delta \overline{h}
-\beta_4\Delta \overline{k}
+\beta_5\Delta \mathrm{unlock}
+\beta_6\Delta \mathrm{terminalReadiness}.
\end{align}
A packet receives no answer-readiness credit merely for adding evidence. Contextual statistics are updated only inside the current question, with prior strength preventing one stochastic outcome from dominating. The trace stores the prior estimate, predicted EVC, raw and normalized components, actual resource vector, realized utility, and posterior estimate.

\section{Join Schemas and Auditing}
\label{app:joins}

A join proposal is accepted only if all premises are active, independently scored, provenance-addressable, and jointly sufficient under a registered schema. The current schema families are:
\begin{itemize}[leftmargin=*]
  \item \textbf{Path composition:} $R_1(x,y)\land R_2(y,z)\Rightarrow R(x,z)$;
  \item \textbf{Shared role:} $R(x,z)\land R(y,z)\Rightarrow C(x,y)$;
  \item \textbf{Set intersection:} $x\in A\land x\in B\Rightarrow x\in A\cap B$;
  \item \textbf{Comparison:} $f(x)=a\land f(y)=b\land a>b\Rightarrow x\succ_f y$;
  \item \textbf{Numeric/temporal composition:} normalized quantities or dates are compared or aggregated with unit checks.
\end{itemize}

An auditable three- or four-hop case must expose the complete ordered premise chain, source spans, intermediate claim texts, join schemas, acceptance scores, revision history, and the final candidate dependency. Merely retrieving all gold passages does not count as chain completion.

\section{Evaluation Definitions}
\label{app:metrics}

\paragraph{Answer metrics.}
EM and token F1 follow each benchmark's official normalization. Candidate presence means that the normalized gold answer appears among active answer candidates before termination. Full-chain completion requires an ordered, provenance-backed route through all required hops. Answer-in-context asks whether the answer string or accepted alias occurs in the retrieved evidence available to the reader. Ordered evidence coverage gives partial credit for correctly ordered supporting evidence.

\paragraph{Reliability metrics.}
Unsupported-answer rate is the fraction of emitted answers that fail the support gate under an independent audit. Selective accuracy is accuracy conditioned on \textsc{Answer}. Coverage is the fraction of all examples answered. Revision precision is the fraction of revision events that remove or repair a truly conflicting state; revision recall is the fraction of injected or independently annotated conflicts that trigger an appropriate revision. Terminal calibration is reported as a confusion matrix over the three terminal outcomes against post-hoc solvability and budget status.

\paragraph{Allocation metrics.}
We report resource shares by subgoal, graph region, packet family, and fidelity; entropy and Gini statistics summarize non-uniformity. Real operation-choice rate excludes forced and bookkeeping actions. Cross-family choice rate measures steps at which distinct feasible families were genuinely compared. EVC quality includes Spearman rank correlation with realized utility, sign accuracy, mean absolute error after calibration, and cumulative regret:
\begin{equation}
\mathrm{Regret}=\sum_t\left(\max_{p\in\mathcal{P}_t^{\mathrm{obs}}} r_t(p)-r_t(p_t)\right),
\end{equation}
where counterfactual utility is evaluated only when a controlled replay or paired intervention makes it observable; otherwise the term is reported as missing rather than imputed.

\paragraph{Compute metrics.}
Each run reports aggregate and per-example calls, input/output/generated tokens where available, retrieval count, wall-clock latency, and estimated API cost. We show median, mean, 90th percentile, and maximum because adaptive methods can hide expensive tails behind similar averages.

\section{Frozen Experimental Protocol}

\subsection{Development and evaluation phases}

The protocol separates:
\begin{enumerate}[leftmargin=*]
  \item \textbf{smoke tests} (up to 20 examples) for schema and trace correctness;
  \item \textbf{development tests} (50 examples) for mechanism gates and bounded iteration;
  \item \textbf{frozen validation} on larger, preselected IDs for model selection;
  \item \textbf{final evaluation} run once per frozen configuration, with no prompt or threshold changes afterward.
\end{enumerate}
All paired comparisons share question IDs. Checkpoints include code commit, configuration hash, corpus/index identifier, model name, API settings, seed, and raw trace directory. API failures are retried according to a fixed policy and reported.

\subsection{Hard gate before final-scale claims}

The final system must satisfy all of the following:
\begin{itemize}[leftmargin=*]
  \item \textbf{Infrastructure:} zero leakage, zero invariant violations, controller-only mutation;
  \item \textbf{Reasoning capability:} candidate presence at least v1 $+0.10$, full-chain completion at least v1 $+0.10$, and at least three auditable three-/four-hop join cases;
  \item \textbf{Dynamic behavior:} adversarial revision works, natural revision precision is acceptable, and allocation is demonstrably non-uniform;
  \item \textbf{Computation allocation:} complete EVC traces, predicted EVC and actual costs recorded, and at least one Pareto-improving operating point;
  \item \textbf{Termination:} no unsupported answers and correct separation of \textsc{Answer}, \textsc{Abstain}, and \textsc{BudgetExhausted}.
\end{itemize}
Thresholds described as acceptable must be numerically fixed before final evaluation. Failure of a gate is reported as a negative result; it is not repaired using test examples.

\section{Additional Results Placeholders}

\begin{figure}[h]
\centering
\placeholderbox{4.0cm}{Per-hop answer quality and cost on MuSiQue, HotpotQA, and 2WikiMultiHopQA, with paired confidence intervals.}
\caption{Per-hop performance placeholder.}
\label{fig:perhop}
\end{figure}

\begin{figure}[h]
\centering
\placeholderbox{4.0cm}{Example dynamic hypergraph trace: diffusion before selection, chosen packet and predicted EVC, accepted join, contradiction event, revision, and terminal gate.}
\caption{Auditable case-study placeholder.}
\label{fig:case}
\end{figure}

\begin{table}[h]
\caption{Termination confusion matrix placeholder. Rows are oracle/post-hoc outcomes and columns are system terminal decisions.}
\centering
\begin{tabular}{lrrr}
\toprule
 & \textsc{Answer} & \textsc{Abstain} & \textsc{BudgetExhausted} \\
\midrule
Supported/solved & \tbd{ } & \tbd{ } & \tbd{ } \\
Unsupported/unsolved & \tbd{ } & \tbd{ } & \tbd{ } \\
Useful work blocked by budget & \tbd{ } & \tbd{ } & \tbd{ } \\
\bottomrule
\end{tabular}
\end{table}

\section{Reproducibility Checklist}

The final artifact should include: exact evaluation IDs; immutable corpus and index fingerprints; all prompts; decoding settings; model and endpoint identifiers; controller and gate configurations; raw operation traces; per-example metrics; budget ledgers; baseline commands; environment lock files; and scripts that regenerate every table and figure from raw traces. Secrets and credentials must never be committed. Because hosted models may change, the final paper should report run dates and, where the provider exposes them, model snapshots.
